\section{总结与延伸}

\subsection{本讲核心要点回顾}

本讲作为 MIT 6.S191 基础讲座系列的总结篇，覆盖了三个核心主题：

\begin{enumerate}
    \item \textbf{深度学习的本质与局限}
    \begin{itemize}
        \item 神经网络是函数逼近器，Universal Approximation Theorem 保证了理论能力但不保证泛化
        \item 随机标签实验揭示了深度网络惊人的过拟合（memorization）能力
        \item 三大失败模式：数据偏差、安全关键场景的不确定性、对抗攻击
        \item 九大局限性：数据饥渴、计算密集、对抗脆弱、算法偏见、不确定性差、黑箱、需要专家知识、难编码结构、外推困难
    \end{itemize}

    \item \textbf{新前沿 I -- Diffusion Models}
    \begin{itemize}
        \item 克服 VAE/GAN 的 mode collapse 和训练不稳定问题
        \item 前向加噪构建训练对，反向去噪实现生成
        \item 迭代生成策略是质量优势的核心
        \item 广泛应用：Text-to-Image、视频生成、分子/蛋白质设计
    \end{itemize}

    \item \textbf{新前沿 II -- Large Language Models}
    \begin{itemize}
        \item 本质：超大规模神经网络 + 超大规模文本数据
        \item 核心训练目标：Next Token Prediction
        \item 涌现能力随模型规模出现``相变''式增长
        \item 关键挑战：鲁棒性、幻觉、安全防护、逻辑推理
    \end{itemize}
\end{enumerate}

\subsection{拓展阅读}

\begin{itemize}
    \item Hornik, K., Stinchcombe, M., \& White, H. (1989). Multilayer feedforward networks are universal approximators. \textit{Neural Networks}.
    \item Zhang, C., et al. (2017). Understanding deep learning requires rethinking generalization. \textit{ICLR 2017}.
    \item Goodfellow, I., Shlens, J., \& Szegedy, C. (2015). Explaining and harnessing adversarial examples. \textit{ICLR 2015}.
    \item Athalye, A., et al. (2018). Synthesizing robust adversarial examples. \textit{ICML 2018}.
    \item Ho, J., Jain, A., \& Abbeel, P. (2020). Denoising diffusion probabilistic models. \textit{NeurIPS 2020}.
    \item Sohl-Dickstein, J., et al. (2015). Deep unsupervised learning using nonequilibrium thermodynamics. \textit{ICML 2015}.
    \item Ramesh, A., et al. (2022). Hierarchical text-conditional image generation with CLIP latents. \textit{arXiv 2022}.
    \item Wei, J., et al. (2022). Emergent abilities of large language models. \textit{TMLR 2022}.
    \item 课程官网: \url{http://introtodeeplearning.com}
    \item GitHub: \url{https://github.com/MITDeepLearning/introtodeeplearning/}
\end{itemize}

%% --- Fin del contenido --- %%

\end{document}
